% STATUS: written 2026-11-05 (Ore 1933 pp. 480--483, Lequain 2008
% pp. 801--802 and Nowicki 2004 pp. 4568--4570 read at the cited pages).
\section{Differential polynomial rings}\label{sec:polynomial-weyl:ore-differential}

Section~\ref{sec:polynomial-weyl:differential-operators} built operators
on a fixed ring. This section moves the derivation into the multiplication:
a polynomial variable acts on the coefficient ring through a derivation.
Ore built the general skew polynomial ring in 1933
\cite[p.~480]{Ore1933}, and
Section~\ref{sec:generalized-twisted:ore-extensions} recalls his
construction. The differential twist of his Section~2 is the case of this
section \cite[p.~482]{Ore1933}. % claim-ok: the cited page states the differential special case
The coefficient ring is commutative, the conjugation is the identity,
and the twist of $a$ by $x$ is the derivative $a'$
\cite[p.~482]{Ore1933}.

\begin{remark}[{\cite[p.~482]{Ore1933}}]\label{rem:pw:ore-differential}
Ore states the differential case as follows. In the theory of linear
differential equations the coefficient field $K$ is commutative, and the
conjugate $\bar{a}$ equals $a$. The derivative then satisfies
$(a+b)'=a'+b'$ and $(ab)'=ab'+a'b$, and it has the ordinary properties
of a derivative \cite[p.~482]{Ore1933}. % claim-ok: the header of the remark already cites this page; the sentence is part of the quote
In the theory of difference equations the derivative is
$0$ and the conjugate is not \cite[p.~482]{Ore1933}.
\end{remark}

The differential Ore extension of a ring $R$ by a derivation $d$ is
denoted $R[X;d]$. Its elements are polynomials in $X$ with coefficients
in $R$. The multiplication rule is $Xa=aX+d(a)$
\cite[p.~481, (3)]{Ore1933}. % claim-ok: the cited formula (3) is x . a = abar . x + a', the differential case of which is the rule stated here
The Weyl algebra $A_{1}$ of Section~\ref{sec:polynomial-weyl:weyl} is the
first example, since $A_{1}$ is the algebra of polynomial differential
operators in one variable \cite[p.~1]{BelovKanelKontsevich2007}. % claim-ok: the cited identification states the example

A ring $R$ with a derivation $d$ is \emph{$d$-simple} when every
$d$-ideal of $R$ is trivial \cite[p.~801]{Lequain2008}. Goodearl and
Warfield call the same extension a formal differential operator ring
$R[x;\delta]$ \cite[p.~28]{GoodearlWarfield2004}. % claim-ok: the cited page defines R[x;delta], and the proposition below carries the located theorem

\begin{proposition}[{\cite[Proposition~2.1, p.~31]{GoodearlWarfield2004}}]
\label{prop:pw:goodearl-warfield}
Let $R$ be a $\mathbb{Q}$-algebra and $\delta$ a derivation on $R$. Then
$R[x;\delta]$ is a simple ring if and only if $R$ is $\delta$-simple and
$\delta$ is outer.
\end{proposition}

Lequain gives the criterion in a form ready for computation, and recalls
the criterion of Shamsuddin for one step further, where the derivation
extends to $R[Y]$ with $d(Y)=aY+b$ \cite[p.~802]{Lequain2008}.

\begin{theorem}[{Shamsuddin, \cite[Theorem~2.1, p.~802]{Lequain2008}}]
\label{thm:pw:shamsuddin}
Let $R$ be a ring, $Y$ an indeterminate over $R$ and $d$ a derivation of
$R[Y]$ such that $d(R)\subseteq R$, $R$ is $d$-simple and
$d(Y)=aY+b$ with $a,b\in R$. Then the following are equivalent:
\begin{enumerate}
\item $R[Y]$ is $d$-simple;
\item the equation $d(Z)=aZ+b$ has no solution in $R$.
\end{enumerate}
\end{theorem}

The concrete ground for this chapter is the Ore extension $A=k[x][t;d]$, with
$d$ a nonzero derivation of $k[x]$. Over a field of characteristic zero,
every local derivation of this ring is a derivation, by
Theorem~\ref{thm:pw:nowicki-ore} below. % claim-ok: the theorem carries the citation


\begin{theorem}[{\cite[Theorem~4.1, p.~4569]{Nowicki2004Local}}]
\label{thm:pw:nowicki-ore}
Let $k[x]$ be the polynomial ring in one variable over a field $k$ of
characteristic zero, and let $d$ be a nonzero derivation of $k[x]$. Let
$A=k[x][t;d]$ be the Ore extension of $k[x]$ with respect to $d$. Every
local derivation of $A$ is a derivation.
\end{theorem}

The local derivations of
Section~\ref{sec:basic-theory:definitions} therefore give nothing new on
this ring. % claim-ok: Theorem thm:pw:nowicki-ore just above is the statement of this sentence
Chapter~\ref{ch:commutative} follows polynomial derivations into the
geometry of the Jacobian question that
Section~\ref{sec:polynomial-weyl:polynomial} left open % claim-ok: own signpost to the next chapter
\cite[p.~1]{BelovKanelKontsevich2007}.
